% AMS 会议文集 / Contemporary Mathematics —— amsproc
% 引擎：pdflatex
\documentclass{amsproc}

\usepackage{amssymb,amsthm,mathtools}
\usepackage{graphicx}

\newtheorem{theorem}{Theorem}[section]
\newtheorem{proposition}[theorem]{Proposition}
\theoremstyle{definition}
\newtheorem{definition}[theorem]{Definition}
\newcommand{\R}{\mathbb{R}}
\newcommand{\Om}{\Omega}

\title[Short running title]{A note on the asymptotic behavior of solutions}
\author{Yang Jianwei}
\address{School of Mathematics and Statistics, NCWU, Zhengzhou 450046, China}
\email{yangjianwei@example.edu.cn}

\begin{document}

\begin{abstract}
We present a short proof of the long-time asymptotic behavior of solutions
to a damped semilinear wave equation. The argument is based on a
Lyapunov functional and the {\L}ojasiewicz--Simon inequality.
\end{abstract}

\maketitle

\section{Introduction}

Let $\Om\subset\R^n$ be a bounded smooth domain. We consider
\begin{equation}\label{eq:wave}
  u_{tt} + u_t - \Delta u + g(u) = 0 , \qquad (x,t)\in\Om\times(0,\infty),
\end{equation}
subject to the Dirichlet boundary condition $u=0$ on $\partial\Om$.

\begin{theorem}
Assume $g\in C^1(\R)$ is real analytic and satisfies the growth condition
$|g'(s)|\le C(1+|s|^{p-1})$ with $p<n/(n-2)$. Then every global bounded
solution of \eqref{eq:wave} converges to a single stationary state as
$t\to\infty$.
\end{theorem}

\section{Proof sketch}

\begin{proof}
The energy functional
\begin{equation}
  E(u,u_t)=\int_{\Om}\Bigl(\tfrac12|u_t|^2+\tfrac12|\nabla u|^2+G(u)\Bigr)\,dx,
  \qquad G(s)=\int_0^s g(r)\,dr,
\end{equation}
is nonincreasing along trajectories. Since $g$ is analytic, the
{\L}ojasiewicz--Simon inequality applies to $E$ near any stationary point,
which yields the convergence.
\end{proof}

\section*{Acknowledgment}
This note was prepared for the conference proceedings volume.

For the classical background and related results, see~\cite{Si}.

\begin{thebibliography}{9}
\bibitem{Si} L.~Simon, \emph{Asymptotics for a class of nonlinear evolution
equations, with applications to geometric problems}, Ann. of Math. \textbf{118} (1983), 525--571.
\end{thebibliography}

\end{document}
